\phantomsection\addcontentsline{toc}{section}{Smart Materials}
\ECchapter{12}{Smart Materials}{Can materials react to their environment?}{ECPurple}

\ECObjectives{
\begin{itemize}
\item Identify the stimulus, physical coupling and response of several smart materials.
\item Explain how shape-memory and piezoelectric effects are integrated into controlled systems.
\item Evaluate stroke, force, response time, fatigue, safety and lifecycle constraints.
\end{itemize}}

\begin{multicols}{2}
\begin{engineeringnews}
Smart materials can sense or respond to temperature, stress, electricity, light or magnetic fields.
They are used in medical devices, vibration control, precision positioning and energy-efficient
buildings, where one compact element can perform structural, sensing and actuation functions.
\end{engineeringnews}

\begin{ecarticle}
A conventional material is selected mainly for passive properties such as stiffness, strength or
thermal conductivity. A smart material produces a predictable response to an external stimulus. The
response may be a change in shape, electrical charge, colour, viscosity or magnetic behaviour.
Engineers use this physical coupling to create sensors, actuators and adaptive structures.

Shape-memory alloys are a well-known example. Nickel-titanium can be deformed in a low-temperature
phase and recover a previously defined shape when heated. A thin wire can therefore act as a compact
actuator: electrical current heats it, the wire contracts, and a spring or external load returns it
when it cools. The available stroke is useful, but cooling may limit the operating frequency.

Piezoelectric materials couple mechanical and electrical energy. When compressed or bent, they
produce electrical charge; this direct effect is used in force, pressure and vibration sensors. When
voltage is applied, they deform slightly; this inverse effect is used in ultrasound transducers, fuel
injectors and nanometre-scale positioning systems. Their motion is small, but rapid and precise.

Other smart materials provide different functions. Electrochromic layers change optical transmission
when a voltage is applied. Magnetorheological fluids become more resistant to flow in a magnetic
field and can adjust suspension damping. Self-healing polymers may release a repair agent when a crack
ruptures embedded microcapsules.

These materials are not intelligent by themselves. A useful engineering system still requires
sensors, power electronics, control logic and safe limits. Performance changes with temperature,
fatigue, hysteresis and ageing. Designers must consider response time, energy consumption,
recyclability and failure behaviour. A passive safe state is often essential.
\end{ecarticle}

\begin{tcolorbox}[title=\sffamily\bfseries Engineering Insight,colback=yellow!10,colframe=ECOrange]
The material law alone does not define actuator performance. The useful output depends on geometry,
preload, heat transfer, drive electronics, feedback control and the mechanical load connected to the
material.
\end{tcolorbox}

\begin{engineeringfocus}
\textbf{Stimulus--material--response chain}
\begin{center}
\begin{tikzpicture}[font=\scriptsize,>=Latex,node distance=3mm,scale=.80,transform shape]
\node[draw=ECBlue,fill=ECLightBlue,rounded corners,align=center] (stim) {stimulus\\temperature, force, voltage};
\node[draw=ECPurple,fill=ECPurple!10,rounded corners,align=center,right=of stim] (mat) {smart\\material};
\node[draw=ECGreen,fill=ECGreen!10,rounded corners,align=center,right=of mat] (resp) {response\\shape, charge, colour};
\node[draw=ECOrange,fill=ECOrange!10,rounded corners,align=center,below=of mat] (control) {controller and\\power interface};
\draw[->,thick] (stim) -- (mat);
\draw[->,thick] (mat) -- (resp);
\draw[<->,thick,ECOrange] (control) -- (mat);
\draw[->,dashed,ECBlue] (resp.south) |- (control.east);
\end{tikzpicture}
\end{center}
Closed-loop control measures the actual response and adjusts the stimulus. This improves accuracy and
compensates for temperature variation, hysteresis and ageing.
\end{engineeringfocus}

\begin{keyfigures}
\begin{tabularx}{\linewidth}{@{}Xr@{}}
Shape-memory alloy strain & several percent\\
Piezoelectric actuator motion & micrometres\\
Piezoelectric response & very fast\\
Electrochromic transition & seconds to minutes\\
Critical design issues & fatigue, hysteresis, ageing
\end{tabularx}
\end{keyfigures}

\begin{tcolorbox}[title=\sffamily\bfseries Shape recovery with temperature,colback=white,colframe=ECPurple]
\begin{center}
\begin{tikzpicture}
\begin{axis}[width=.96\linewidth,height=45mm,xlabel={Temperature ($^\circ$C)},ylabel={Recovered strain (\%)},xmin=15,xmax=95,ymin=0,ymax=5.5,xticklabel style={font=\scriptsize},yticklabel style={font=\scriptsize},grid=major]
\addplot[very thick,mark=*] table[x=temperature,y=strain,col sep=comma]{data/EC12/smart_response.csv};
\end{axis}
\end{tikzpicture}
\end{center}
\footnotesize Illustrative response of a shape-memory actuator. The transformation occurs over a
temperature range rather than at one exact temperature.
\end{tcolorbox}

\begin{tcolorbox}[title=\sffamily\bfseries Piezoelectric vibration control,colback=white,colframe=ECBlue]
\begin{center}
\begin{tikzpicture}[font=\scriptsize,>=Latex]
\draw[very thick] (0,0) -- (5.2,0);
\draw[fill=ECBlue!12,draw=ECBlue,thick] (.3,0) rectangle (4.9,.22);
\draw[fill=ECPurple!20,draw=ECPurple,thick] (1.3,.22) rectangle (3.9,.42);
\node[above] at (2.6,.45) {bonded piezoelectric patch};
\draw[->,thick,ECRed] (5.05,.1) -- +(0,1.05) node[above]{disturbance};
\draw[->,thick,ECGreen] (2.6,1.65) -- (2.6,.48) node[midway,right]{controlled voltage};
\draw[decorate,decoration={snake,amplitude=1.2mm,segment length=5mm},ECBlue,thick] (.5,-.35) -- (4.7,-.35);
\node[below,align=center] at (2.6,-.55) {measured vibration reduced\\by an opposing actuator force};
\end{tikzpicture}
\end{center}
A sensor measures vibration, the controller calculates a corrective command and the patch produces a
small deformation that opposes the motion.
\end{tcolorbox}

\begin{vocabulary}
\begin{tabularx}{\linewidth}{@{}lX@{}}
smart material & matériau intelligent\\
stimulus & sollicitation / stimulus\\
shape-memory alloy & alliage à mémoire de forme\\
phase transformation & transformation de phase\\
piezoelectric effect & effet piézoélectrique\\
strain & déformation relative\\
damping & amortissement\\
hysteresis & hystérésis\\
ageing & vieillissement\\
fail-safe state & état sûr en cas de panne
\end{tabularx}
\end{vocabulary}

\begin{questions}
\item What distinguishes a smart material from a passive material?
\item How can a shape-memory wire act as an actuator?
\item Explain the direct and inverse piezoelectric effects.
\item Why is closed-loop control useful with smart materials?
\item Use the graph to estimate the useful transformation temperature range.
\item Why must fatigue, hysteresis and failure behaviour be considered together?
\end{questions}

\begin{applicationexercise}
A piezoelectric actuator is 40 mm long and produces a free strain of 0.08\% when powered.
\begin{enumerate}
\item Calculate its free displacement in millimetres and micrometres.
\item A mechanical amplifier multiplies the displacement by 4. Estimate the output movement.
\item If the loaded displacement is only 65\% of the free value, calculate the useful movement.
\item Explain why amplification usually reduces available force or stiffness.
\end{enumerate}
\end{applicationexercise}

\begin{engineeringchallenge}
Design an adaptive bicycle-helmet ventilation system using a smart material. Define requirements,
stimulus, response, actuator geometry and control strategy. Estimate the required movement, compare at
least two candidate materials, identify fatigue and thermal risks, and explain how the helmet remains
safe if the active material or power supply fails.
\end{engineeringchallenge}

\begin{speaking}
Choose one smart material and present an everyday product that could use it. Explain the benefit,
technical limitation, control requirement and environmental impact in a four-minute product review.
\end{speaking}
\end{multicols}

\subsection*{Sources}
\ECsource{NIST -- Materials Research}{https://www.nist.gov/topics/materials}
\ECsource{European Space Agency -- Space Engineering and Technology}{https://www.esa.int/Enabling_Support/Space_Engineering_Technology}
